% !PRINTED 5
\weberpage{5}
%<<CONT
setzen, und darin ist $hmn'>hm'n$. Man kann $h$ immer so
annehmen, dass zwischen $hmn'$ und $hm'n$ noch Zahlen liegen,
und wenn $p$ eine solche Zahl ist, so liegt $p:hnn'$ zwischen
$\mu$ und $\mu'$.

Die Punkte einer geraden Linie kann man auch als eine geordnete
Menge auffassen, wenn man unter grösser und kleiner irgend eine
Ortsbezeichnung, z.~B. weiter rechts und weiter links oder höher
und tiefer versteht.

\begin{center}
\rule{0.18\linewidth}{0.4pt}
\end{center}

Eine Eintheilung einer geordneten Menge $\mathfrak M$ in zwei Theile
$A,B$ der Art, dass jedes Element $a$ von $A$ kleiner ist als jedes
Element $b$ von $B$, wird ein \gesperrt{Schnitt} in $\mathfrak M$
genannt und wird passend durch $(A,B)$ bezeichnet. Ein solcher
Schnitt entsteht, wenn man irgend ein Element $\mu$ von $\mathfrak M$
herausgreift, alle kleineren Elemente zu $A$, alle grösseren zu $B$
und $\mu$ selbst nach Belieben zu $A$ oder zu $B$ rechnet. Es
entstehen, genau gesagt, je nachdem man das eine oder das andere
thut, zwei Schnitte, die wir aber immer als gleich betrachten wollen.
Wenn in einem Schnitt $(A,B)$ entweder $A$ ein grösstes oder $B$ ein
kleinstes Element $\mu$ enthält, so sagen wir, dass $\mu$ den Schnitt
$(A,B)$ erzeugt. Es kann aber auch der Fall vorkommen, dass weder $A$
ein grösstes noch $B$ ein kleinstes Element besitzt.

Wenn jeder Schnitt in einer dichten Menge durch ein bestimmtes
Element $\mu$ erzeugt wird, so heisst die Menge \gesperrt{stetig}.

Stetigkeit sowohl als Dichtigkeit sind Eigenschaften, die der Natur
der Sache nach unserer Sinneswahrnehmung unzugänglich sind; sie lassen
sich daher auch an Dingen der Aussenwelt, an Raumgrössen, Zeiträumen,
Massen niemals mit Strenge nachweisen, wie sehr sie uns auch im Wesen
unserer Anschauung zu liegen scheinen. Es lassen sich aber sehr wohl
reine Begriffssysteme construiren, denen die Dichtigkeit ohne die
Stetigkeit oder auch Dichtigkeit und Stetigkeit zukommen\footnotemark[1].
\footnotetext[1]{Dies ist von \gesperrt{Dedekind} nachgewiesen, dem wir
überhaupt die oben gegebene Definition der Stetigkeit verdanken. Vgl.
die Schriften von \gesperrt{Dedekind}, „Stetigkeit und irrationale
Zahlen“, Braunschweig 1872, 1892. „Was sind und was sollen die
Zahlen?“, Braunschweig 1888, 1893. Andere}

Ein Beispiel einer dichten Menge bieten die \gesperrt{rationalen
Brüche}. Diese Mannigfaltigkeit, die wir mit $\mathfrak R$ bezeichnen
%CONT>>
